\documentclass[11pt]{article}

\usepackage{amsmath,amssymb,mathtools}
\usepackage[margin=1in]{geometry}
\usepackage[hidelinks]{hyperref}

\title{What the code computes: Virasoro sewing vs.\ Nekrasov/AGT (to level 2)}
\author{}
\date{}

\begin{document}
\maketitle

\section*{1.\ Setup and normalization}

Fix central charge $c$. Consider the Virasoro 4-point conformal block on $\mathbb{P}^1$
with primary insertions at
\[
(0,q,1,\infty)
\]
of conformal dimensions
\[
(h_2,h_1,h_3,h_4),
\]
and intermediate (internal) highest weight of dimension $h$ in the $s$-channel.
The code computes the \emph{reduced} block
\[
\mathcal{F}(q)= 1+\sum_{n\ge 1} c_n\,q^n,
\]
i.e.\ after removing the overall prefactor $q^{\,h-h_1-h_2}$ (and using the convention $c_0=1$).

\section*{2.\ Virasoro side: sewing with the Shapovalov form}

Let $M_h$ be the Verma module with highest weight vector $\lvert h\rangle$.
At level $n$ we use the PBW basis labelled by partitions $\lambda\vdash n$:
\[
L_{-\lambda}\lvert h\rangle:=L_{-\lambda_1}L_{-\lambda_2}\cdots L_{-\lambda_\ell}\lvert h\rangle,
\qquad \lambda_1\ge\lambda_2\ge\cdots\ge\lambda_\ell\ge 1.
\]

\subsection*{2.1 Gram (Shapovalov) matrix}
The level-$n$ Gram matrix is
\[
G^{(n)}_{\lambda\mu} := \langle h \rvert \, L_{\lambda}\,L_{-\mu}\,\lvert h\rangle,
\qquad
L_\lambda := L_{\lambda_\ell}\cdots L_{\lambda_1},
\]
computed using
\[
[L_m,L_n]=(m-n)L_{m+n}+\frac{c}{12}m(m^2-1)\delta_{m+n,0},
\qquad
L_0\lvert h\rangle = h\lvert h\rangle,
\qquad
\langle h\rvert h\rangle=1,
\]
and the code inverts $G^{(n)}$.

\subsection*{2.2 Vertex vectors from Ward identities}
For the left trivalent vertex define
\[
F_u(z):=\frac{\langle h_2\rvert\,V_{h_1}(z)\,u\rangle}{\langle h_2\rvert\,V_{h_1}(z)\,\lvert h\rangle},
\qquad u\in M_h,
\]
so $F_{\lvert h\rangle}(z)=1$. Since $\langle h_2\rvert L_{-n}=0$ for $n>0$, the Ward identity
\[
[L_m,V_{h_1}(z)] = z^m\bigl(z\partial_z+(m+1)h_1\bigr)V_{h_1}(z)
\]
implies, for $n>0$,
\[
F_{L_{-n}u}(z) = -z^{-n}\bigl(z\partial_z+(1-n)h_1\bigr)F_u(z).
\]
Using the primary scaling $\langle h_2\rvert V_{h_1}(z)\lvert h\rangle\propto z^{h_2-h_1-h}$, the recursion
reduces to multiplying by linear factors (the code evaluates at $z=1$):
\[
v^{(12)}_\lambda := F_{L_{-\lambda}\lvert h\rangle}(1).
\]
Similarly, for the right vertex
\[
v^{(34)}_\lambda := \frac{\langle h_4\rvert\,V_{h_3}(1)\,L_{-\lambda}\lvert h\rangle}{\langle h_4\rvert\,V_{h_3}(1)\,\lvert h\rangle}.
\]

\subsection*{2.3 Sewing formula}
At each level $n$,
\[
c_n = \sum_{\lambda,\mu\vdash n} v^{(12)}_\lambda \,(G^{(n)})^{-1}_{\lambda\mu}\, v^{(34)}_\mu
\;=\;
\bigl(v^{(12)}\bigr)^{T}(G^{(n)})^{-1}v^{(34)}.
\]

\subsection*{2.4 Explicit answers at levels $0,1,2$ (Virasoro)}
\paragraph{Level 0.} $c_0=1$.

\paragraph{Level 1.}
The basis is $\{L_{-1}\lvert h\rangle\}$,
\[
G^{(1)} = (2h),
\qquad
v^{(12)}_{[1]} = h+h_1-h_2,
\qquad
v^{(34)}_{[1]} = h+h_3-h_4,
\]
so
\[
c_1 = \frac{(h+h_1-h_2)(h+h_3-h_4)}{2h}.
\]

\paragraph{Level 2.}
Use the basis $\{L_{-2}\lvert h\rangle,\; L_{-1}^2\lvert h\rangle\}$.
The Gram matrix is
\[
G^{(2)}=
\begin{pmatrix}
\frac{c}{2}+4h & 6h\\
6h & 4h(2h+1)
\end{pmatrix},
\qquad
\det G^{(2)} = \Bigl(\frac{c}{2}+4h\Bigr)\,4h(2h+1)-(6h)^2.
\]
The vertex components are
\[
v^{(12)}_{[2]}=h+2h_1-h_2,
\qquad
v^{(12)}_{[11]}=(h+h_1-h_2)(h+h_1-h_2+1),
\]
\[
v^{(34)}_{[2]}=h+2h_3-h_4,
\qquad
v^{(34)}_{[11]}=(h+h_3-h_4)(h+h_3-h_4+1).
\]
Writing $A:=h+h_1-h_2$, $B:=h+h_3-h_4$, $A_2:=h+2h_1-h_2$, $B_2:=h+2h_3-h_4$, one finds
\[
c_2
=
\frac{1}{\det G^{(2)}}
\Bigl(
4h(2h+1)\,A_2B_2
-6h\,A_2\,B(B+1)
-6h\,A(A+1)\,B_2
+\Bigl(\frac{c}{2}+4h\Bigr)A(A+1)B(B+1)
\Bigr).
\]

\section*{3.\ Gauge theory side: Nekrasov instanton sum and AGT dictionary}

Fix equivariant parameters $\epsilon_1,\epsilon_2$ and set $Q:=\epsilon_1+\epsilon_2$.
For a Young diagram $Y$, write boxes as $s=(i,j)$ with arm/leg lengths $A_Y(s),L_Y(s)$.
Define
\begin{align*}
N_{Y,W}(x)
&:=
\prod_{s\in Y}\Bigl(x+\epsilon_1(A_Y(s)+1)-\epsilon_2L_W(s)\Bigr)
\prod_{t\in W}\Bigl(x-\epsilon_1A_W(t)+\epsilon_2(L_Y(t)+1)\Bigr),\\
Z_{\mathrm{vec}}(Y_1,Y_2)
&:=
\prod_{A,B=1}^2 \frac{1}{N_{Y_A,Y_B}(a_A-a_B)},
\qquad (a_1,a_2)=(a,-a),\\
Z_{\mathrm{fund}}(Y_1,Y_2;\mu)
&:=
\prod_{A=1}^2\ \prod_{(i,j)\in Y_A}\Bigl(a_A+\mu+\epsilon_1(j-1)+\epsilon_2(i-1)\Bigr).
\end{align*}
Then for four masses $\mu_1,\dots,\mu_4$ the instanton series is
\[
Z_{\mathrm{inst}}(q)
=
\sum_{Y_1,Y_2} q^{|Y_1|+|Y_2|}
\;Z_{\mathrm{vec}}(Y_1,Y_2)\;\prod_{f=1}^4 Z_{\mathrm{fund}}(Y_1,Y_2;\mu_f)
=
\sum_{k\ge 0} Z_k\,q^k.
\]

\subsection*{3.1 AGT map used by the code}
The code uses
\[
c = 1+\frac{6Q^2}{\epsilon_1\epsilon_2},
\qquad
h(\alpha)=\frac{\alpha(Q-\alpha)}{\epsilon_1\epsilon_2}.
\]
External ``momenta'' $(\beta_0,\alpha_0,\alpha_1,\beta_2)$ are assigned to $(0,q,1,\infty)$:
\[
h_2=h(\beta_0),\quad h_1=h(\alpha_0),\quad h_3=h(\alpha_1),\quad h_4=h(\beta_2).
\]
The internal momentum is
\[
\beta_1=a+\frac{Q}{2},\qquad h=h(\beta_1).
\]
Masses and the $U(1)$ exponent are
\[
\mu_1=-\frac{Q}{2}+\alpha_0+\beta_0,\quad
\mu_2=\phantom{-}\frac{Q}{2}+\alpha_0-\beta_0,\quad
\mu_3=\frac{3Q}{2}-\alpha_1-\beta_2,\quad
\mu_4=\phantom{-}\frac{Q}{2}-\alpha_1+\beta_2,
\]
\[
\nu=\frac{2\alpha_0(Q-\alpha_1)}{\epsilon_1\epsilon_2}.
\]
Finally, the AGT-side reduced series is
\[
\mathcal{F}_{\mathrm{AGT}}(q):=(1-q)^{-\nu}\,Z_{\mathrm{inst}}(q)
=
1+\sum_{n\ge1} c^{\mathrm{AGT}}_n q^n.
\]

\subsection*{3.2 Explicit answers at levels $0,1,2$ (AGT form)}
Expanding $(1-q)^{-\nu}=\sum_{m\ge0}\frac{(\nu)_m}{m!}q^m$ gives:
\[
c^{\mathrm{AGT}}_0=1,\qquad
c^{\mathrm{AGT}}_1=Z_1+\nu,\qquad
c^{\mathrm{AGT}}_2=Z_2+\nu Z_1+\frac{\nu(\nu+1)}{2}.
\]
At one-instanton level, only $(Y_1,Y_2)=([1],\varnothing)$ and $(\varnothing,[1])$ contribute, and the code’s conventions give
\[
Z_1
=
\frac{\prod_{f=1}^4(a+\mu_f)}{\epsilon_1\epsilon_2\,(2a)\,(2a+Q)}
+
\frac{\prod_{f=1}^4(-a+\mu_f)}{\epsilon_1\epsilon_2\,(2a)\,(2a-Q)}.
\]
At two-instanton level, $Z_2$ is computed by summing the five pairs with $|Y_1|+|Y_2|=2$:
\[
([2],\varnothing),\ ([1,1],\varnothing),\ ([1],[1]),\ (\varnothing,[2]),\ (\varnothing,[1,1]),
\]
each weighted by $Z_{\mathrm{vec}}\prod_f Z_{\mathrm{fund}}$ as above.

\section*{4.\ What the code checks}
Given parameters, the code computes $\{c_n\}$ from sewing (Section 2) and independently computes
$\{c^{\mathrm{AGT}}_n\}$ from $\mathcal{F}_{\mathrm{AGT}}(q)$ (Section 3), and compares them through the requested level,
skipping singular loci where the relevant denominators vanish.

\end{document}
